\documentclass[12pt]{article}
\usepackage{amsmath, amssymb, geometry}
\geometry{letterpaper, margin=1in}% 1-inch margins


\title{MATH 374 Assignment Solutions}
\begin{document}
\maketitle
\section*{Prove the following statements:}
\subsubsection*{2.1.16: An odd integer minus an even integer is odd.}
Let \( a = 2m + 1 \) and \( b = 2n \) be an odd and an even integer, respectively. Then, \( a - b = 2m - 2n + 1 = 2(m - n) + 1 \), which is odd. \(\Box\)

\subsubsection*{2.1.17: The product of any two consecutive integers is even.}
Let \(n\) be an integer. Then \(n(n+1)\) is the product of two consecutive integers. \(n(n+1) = n^2 + n\). Either \(n\) or \(n+1\) is even. Therefore, \(n(n+1)\) is even. \(\Box\)

\subsubsection*{2.1.21: If a number \( x \) is positive, so is \( x + 1 \) (do a proof by contraposition).}
To prove this, we will use the method of proof by contraposition.
First, let's express the original statement in the form "If \( P \), then \( Q \)".
\newline
- \( P \): \( x \) is a positive number.
\newline
- \( Q \): \( x + 1 \) is a positive number.
\newline
The contrapositive of a statement "If \( P \), then \( Q \)" is "If \( \lnot Q \), then \( \lnot P \)".
Therefore, the contrapositive of our original statement is:
\newline
"If \( x + 1 \) is not a positive number, then \( x \) is not a positive number."
\newline
This can also be written as:
"If \( x + 1 \leq 0 \), then \( x \leq 0 \)."
\newline
We assume the contrapositive hypothesis is true, i.e., \( x + 1 \leq 0 \).
\newline
Starting with \( x + 1 \leq 0 \), by definition of inequality,
\[ x \leq -1 \]
\[ x \leq 0 \]
Thus, we have shown that \( x \leq 0 \) when \( x + 1 \leq 0 \).
\newline
This confirms the truth of the contrapositive statement, and by extension, the original statement as well.
\newline
Hence, the original statement, "If a number \( x \) is positive, so is \( x + 1 \)," is proven to be true. \(\Box\)

\subsubsection*{2.1.29: If \( n, m, \) and \( p \) are integers and \( n|m \) and \( m|p \), then \( n|p \).}
By definition, \(m = na\) and \(p = mb\) for some integers \(a\) and \(b\). Thus \(p = n(ab)\), implying \(n|p\). \(\Box\)

\subsubsection*{2.1.31: The sum of three consecutive integers is divisible by 3.}
Let \(n\) be an integer. Then \(n + (n + 1) + (n + 2) = 3n + 3 = 3(n + 1)\) is divisible by 3. \(\Box\)

\subsubsection*{2.1.36: For any two numbers \( x \) and \( y \), \( |xy| = |x||y| \).}
Either both \(x\) and \(y\) are positive or negative, or one of them is. In both cases \( |xy| = |x||y| \). \(\Box\)

\subsubsection*{2.1.41: Prove that \( \sqrt{5} \) is not a rational number.}
Assume \(\sqrt{5} = \frac{a}{b}\) with \(a\) and \(b\) having no common factors other than 1. Then \(5b^2 = a^2\), making \(a^2\) a multiple of 5. This contradicts the assumption that \(a\) and \(b\) have no common factors other than 1. \(\Box\)



\newpage
\subsection*{Prove or disprove the following statements:}

\subsubsection*{2.1.45: 297 is a composite number.}
297 can be factored as \(3 \times 99\), and since \(3\) and \(99\) are both integers greater than 1, 297 is a composite number. \(\Box\)

\subsubsection*{2.1.46: 83 is a composite number.}
83 is only divisible by 1 and itself, making it a prime number. Therefore, the statement that 83 is a composite number is disproved. \(\Box\)

\subsubsection*{2.1.63: The sum of two rational numbers is rational.}
Let \(a/b\) and \(c/d\) be two rational numbers. Then, \((a/b) + (c/d) = (ad+bc) / bd\). Since \(ad+bc\) and \(bd\) are integers and \(bd \neq 0\), the sum is rational. \(\Box\)

\subsubsection*{2.1.64: The product of two irrational numbers is irrational.}
This statement is false. Counterexample: \(\sqrt{2} \times \sqrt{2} = 2\), which is rational. Therefore, the product of two irrational numbers can be rational. \(\Box\)

\subsubsection*{2.1.65: The sum of a rational number and an irrational number is irrational.}
Let \(a/b\) be a rational number and \(x\) be an irrational number. Then \(a/b + x\) cannot be expressed as \(m/n\), where \(m\) and \(n\) are integers and \(n \neq 0\), since that would imply \(x\) is rational, contradicting the assumption. Therefore, the sum is irrational. \(\Box\)
.
\end{document}
